\IfFileExists{papers/formalized-vs-literature-preamble.tex}{\documentclass[11pt]{article}
\usepackage[utf8]{inputenc}
\usepackage[T1]{fontenc}
\usepackage{amsmath,amssymb,mathtools}
\usepackage{booktabs,tabularx,array}
\usepackage{enumitem}
\usepackage[margin=1in]{geometry}
\usepackage{xcolor}
\usepackage[hidelinks]{hyperref}
\usepackage{microtype}
\setlength{\parindent}{0pt}
\setlength{\parskip}{0.55em}
\newcommand{\Lean}[1]{\texttt{\nolinkurl{#1}}}
\newcommand{\RepoFile}[1]{\texttt{\nolinkurl{#1}}}
\newcommand{\Class}[1]{\textbf{#1}}
\newcommand{\Top}{\mathord{\top}}
\newcommand{\R}{\mathbb{R}}
\newcommand{\ENN}{\mathbb{R}_{\ge 0}^{\infty}}
\newcommand{\KL}{\operatorname{KL}}
\newcommand{\W}{\mathsf{W}}
\newcommand{\statusline}[2]{%
  \begin{center}\fbox{\begin{minipage}{0.92\linewidth}\textbf{Completion class:} #1\\[2pt]
  \textbf{Strongest caveat:} #2\end{minipage}}\end{center}}
\newenvironment{comparison}{\tabularx{\linewidth}{>{\bfseries}p{0.25\linewidth}X}}{\endtabularx}
\newcommand{\snapshot}{\texttt{902f51aa01a737af68f6feb71c83263cc9b1f4d9}}
}{\documentclass[11pt]{article}
\usepackage[utf8]{inputenc}
\usepackage[T1]{fontenc}
\usepackage{amsmath,amssymb,mathtools}
\usepackage{booktabs,tabularx,array}
\usepackage{enumitem}
\usepackage[margin=1in]{geometry}
\usepackage{xcolor}
\usepackage[hidelinks]{hyperref}
\usepackage{microtype}
\setlength{\parindent}{0pt}
\setlength{\parskip}{0.55em}
\newcommand{\Lean}[1]{\texttt{\nolinkurl{#1}}}
\newcommand{\RepoFile}[1]{\texttt{\nolinkurl{#1}}}
\newcommand{\Class}[1]{\textbf{#1}}
\newcommand{\Top}{\mathord{\top}}
\newcommand{\R}{\mathbb{R}}
\newcommand{\ENN}{\mathbb{R}_{\ge 0}^{\infty}}
\newcommand{\KL}{\operatorname{KL}}
\newcommand{\W}{\mathsf{W}}
\newcommand{\statusline}[2]{%
  \begin{center}\fbox{\begin{minipage}{0.92\linewidth}\textbf{Completion class:} #1\\[2pt]
  \textbf{Strongest caveat:} #2\end{minipage}}\end{center}}
\newenvironment{comparison}{\tabularx{\linewidth}{>{\bfseries}p{0.25\linewidth}X}}{\endtabularx}
\newcommand{\snapshot}{\texttt{902f51aa01a737af68f6feb71c83263cc9b1f4d9}}
}
\title{Kolmogorov extension and Wiener reference measure: formalized result versus the literature}
\author{}
\date{Formalization snapshot \snapshot}
\begin{document}
\maketitle
\statusline{\Class{G/V}: concrete Wiener measure on the unrestricted coordinate path space built from a vendored extension theorem}{The measure lives on all functions $\R_{\ge0}\to\R$; concentration on a canonical continuous anchored path space and its Borel generation are not yet formalized.}
\section{Result}
\Lean{ForMathlib.MeasureTheory.wienerMeasure} in
\RepoFile{ForMathlib/MeasureTheory/WienerMeasure.lean} applies a Kolmogorov extension theorem to Mathlib's consistent Brownian Gaussian projective family.  The repository proves that it is a probability measure and that every finite coordinate restriction has the expected centered Gaussian law with covariance $\min(s,t)$.

\section{Provenance}
The projective-limit implementation under \RepoFile{ForMathlib/KolmogorovExtension/} is vendored from Remy Degenne's
\url{https://github.com/RemyDegenne/kolmogorov_extension4} under Apache-2.0; exact commit metadata is recorded in the vendored headers/README.  The assembly using Mathlib's Brownian projective family is local.

\section{Literature comparison}
Kolmogorov's extension theorem produces a measure on a coordinate product.  The classical Wiener construction then proves continuity of sample paths (or constructs directly on a continuous-path carrier after modification).  The current Lean result completes the projective-limit step but not the continuity/support step.

\begin{comparison}
Index set & Nonnegative real times.\\
Carrier & Unrestricted functions $\R_{\ge0}\to\R$.\\
Finite-dimensional laws & Fully proved.\\
Probability mass & Fully proved.\\
Continuous-path concentration & Open.\\
Anchored path type & Open/provisional.\\
Dyadic Borel generation & Open.\\
\end{comparison}

\section{Role in DRSB}
This gives a concrete reference law for finite-dimensional projection arguments.  It does not by itself support path-level Cameron--Martin translation, disintegration into Brownian bridges, SDE laws, or Girsanov.
\end{document}
